\newif\ifglobal

%%%%%%%%%%%%%%%%%%%%%%%
%%% Compile to view %%%
%%%%%%%%%%%%%%%%%%%%%%%

%%%%%%%%%%%%%%%%%%%%%%%%%%%%%%%%%%%%%%%%%%%%%%%%%%%
%%% Readme:                                     %%%
%%% https://www.overleaf.com/read/bwdjvfjksvjy/ %%%
%%%%%%%%%%%%%%%%%%%%%%%%%%%%%%%%%%%%%%%%%%%%%%%%%%%

% >>> M?: marks a module setting number or important notice

% >>> M4: Set {./relative-path} to external files for this module <<<
\def \modulefiles {./25-RSA}
% To add an external file, use:
% (Replace '---' with actual filename)
% '\input{\modulefiles/tables/---.tex}' for tables
% '\includegraphics{\modulefiles/figures/---}' for figures
% '\subfile{\modulefiles/---}' for register files

% >>> M1: This module is in global project? <<<
% 'YES' - uncomment; 'NO' - comment out
\globaltrue

\ifglobal
    \documentclass[main\_\_EN.tex]{subfiles}
   
\else
    \documentclass[openany, 10pt]{book}

    % >>> M2: In ./00-shared/preamble-trm-module, also find
    %         settings for visibility of todo notes, watermark <<<
    \usepackage{preamble-trm-module}

    % >>> M3: Temporary labels in English,
    %         you can add your own labels there <<<
    \myexternaldocument{./00-shared/temp-labels-en}

\fi %\ifglobal


\setcounter{secnumdepth}{4}


% Define formatting of parts and chapters in text
\titleformat{\chapter}
  [display]
  {\LARGE\bfseries}
  {\color{AllportsColor}Chapter \thechapter}
  {1em}
  {\color{AllportsColor}}
\titlespacing*{\chapter}{0pt}{0pt}{20pt}

% Define headers
\usepackage{fancyhdr}
\usepackage{lipsum}
\renewcommand{\chaptermark}[1]{\markboth{Chapter\, \thechapter\, \ #1}{}}
\fancypagestyle{plain}{%
  \fancyhead{}%
  %\fancyhead[R]{\Navigation}
  \fancyhead[L]{%
    \nouppercase{%
      \ %
      \normalsize\leftmark
    }%
  }%
}
\pagestyle{plain}


% >>> M5: If this module needs any updates in
% './00-shared/preamble-trm-module.sty'
% (include additional package, etc.), contact your module PM
% or create the file './00-shared/changelog.tex' make the
% necessary updates yourself and document them in this file <<<

% Make text left-aligned and allow latex to break pages naturally, comment out for full justification
\raggedright
\raggedbottom


\begin{document}


% Sets language into which pre-defined bits of text are translated
\selectlanguage{english}

% Do not delete this block
\ifglobal
% Add the following front matter if
% this module is outside of global project
\else
    \InputIfFileExists{readme}{}{}
    \listoftodos
    {\let\clearpage\relax \listofdonetodos}
    \tableofcontents
    %\listoftables
    %\listoffigures
\fi


%%%%%%%%%%%%%%%%%%%%%%%%%
%%% TRM Module Begins %%%
%%%%%%%%%%%%%%%%%%%%%%%%%

\hypertarget{rsa}{}
\chapter{RSA Accelerator (RSA)}\label{mod:rsa}
%%%%%%%%%%%%%%%%%%%%%%%%%%%%%%%%%%%%%%%%%%%%%%%%%%%%%%%%%%%%%%%%%%%%%%%%%%%%%%%%%
%%%%%%%%%%%%%%%%%%%%%%%%%%%%%%%%%%%%%%%%%%%%%%%%%%%%%%%%%%%%%%%%%%%%%%%%%%%%%%%%%
\section{Introduction}

The RSA accelerator provides hardware support for high-precision computation used in various RSA asymmetric cipher algorithms, significantly reducing their run time and reducing their software complexity. Compared with RSA algorithms implemented solely in software, this hardware accelerator can speed up RSA algorithms significantly. The RSA accelerator also supports operands of different lengths, which provides more flexibility during the computation.
%%%%%%%%%%%%%%%%%%%%%%%%%%%%%%%%%%%%%%%%%%%%%%%%%%%%%%%%%%%%%%%%%%%%%%%%%%%%%%%%%
%%%%%%%%%%%%%%%%%%%%%%%%%%%%%%%%%%%%%%%%%%%%%%%%%%%%%%%%%%%%%%%%%%%%%%%%%%%%%%%%%
\section{Features}

The following functionality is supported:

\begin{itemize}
    \item Large-number modular exponentiation with two acceleration options
    \item Large-number modular multiplication
    \item Large-number multiplication
    \item Operands of different lengths
    \item Interrupt on completion of computation
\end{itemize}
%%%%%%%%%%%%%%%%%%%%%%%%%%%%%%%%%%%%%%%%%%%%%%%%%%%%%%%%%%%%%%%%%%%%%%%%%%%%%%%%%
%%%%%%%%%%%%%%%%%%%%%%%%%%%%%%%%%%%%%%%%%%%%%%%%%%%%%%%%%%%%%%%%%%%%%%%%%%%%%%%%%

\section{Functional Description}\label{sec:rsa-func-desc}

The RSA accelerator is activated by setting the \hyperref[fielddesc:PCRRSACLKEN]{PCR\_RSA\_CLK\_EN} bit and clearing the \hyperref[fielddesc:PCRRSARSTEN]{PCR\_RSA\_RST\_EN} bit in the \hyperref[regdesc:PCRRSACONFREG]{PCR\_RSA\_CONF\_REG} register. Additionally, users also need to clear \hyperref[fielddesc:PCRDSRSTEN]{PCR\_DS\_RST\_EN} and \hyperref[fielddesc:PCRECDSARSTEN]{PCR\_ECDSA\_RST\_EN} bits to reset \nameref{mod:digsig} and \nameref{mod:ecdsa}.

The RSA accelerator is only available after the \hyperref[sec:mem-sum]{RSA-related memories} are initialized. The content of the \hyperref[regdesc:RSAQUERYCLEANREG]{RSA\_QUERY\_CLEAN\_REG} register is 0 during initialization and will become 1 after the initialization is done. Therefore, wait until \hyperref[regdesc:RSAQUERYCLEANREG]{RSA\_QUERY\_CLEAN\_REG} becomes 1 before using the RSA accelerator.

The \hyperref[regdesc:RSAINTENAREG]{RSA\_INT\_ENA\_REG} register is used to control the interrupt triggered on completion of computation. Write 1 or 0 to this field to enable or disable the interrupt. By default, the interrupt function of the RSA accelerator is enabled.


\begin{importantlisting}
\vspace{-0.5em}

    \chipname{}'s \nameref{mod:digsig} module also calls the RSA accelerator when working. Therefore, users cannot access the RSA accelerator when the \nameref{mod:digsig} module is working.
\end{importantlisting}

%%%%%%%%%%%%%%%%%%%%%%%%%%%%%%%%%%%%%%%%%%%%%%%%%
%%%%%%%%%%%%%%%%%%%%%%%%%%%%%%%%%%%%%%%%%%%%%%%%%
\subsection{Large-Number Modular Exponentiation} \label{sec:rsa-large-modular-Exponentiation}

Large-number modular exponentiation performs \begin{math}Z = X^{Y} \bmod M\end{math}. The computation is based on Montgomery multiplication. Therefore, aside from the \begin{math}X\end{math},
\begin{math}Y\end{math}, and
\begin{math}M\end{math} arguments, two additional ones are needed --- \begin{math}\overline{r}\end{math} and
\begin{math}M^{\prime}\end{math}, which need to be calculated in advance by software.

The RSA accelerator supports operands of length \begin{math}N = 32 \times x\end{math}, where \begin{math}x \in \{ 1,2,3, \dots, 96\} \end{math}. The bit lengths of arguments \begin{math}Z\end{math}, \begin{math}X\end{math}, \begin{math}Y\end{math}, \begin{math}M\end{math}, and \begin{math}\overline{r}\end{math}
can be arbitrary \begin{math}N\end{math}, but all numbers in a calculation must be of the same length. The bit length of \begin{math}M^{\prime}\end{math} must be 32.

To represent the numbers used as operands, let us define a base-\begin{math}b\end{math} positional notation, as follows:
\begin{gather*}
    b = 2^{32}
\end{gather*}

Using this notation, each number is represented by a sequence of base-\begin{math}b\end{math} digits:
\begin{gather*}
    n = \frac{N}{32} \\
    Z = ( Z_{n-1} Z_{n-2} \cdots Z_{0} )_{b} \\
    X = ( X_{n-1} X_{n-2} \cdots X_{0} )_{b} \\
    Y = ( Y_{n-1} Y_{n-2} \cdots Y_{0} )_{b} \\
    M = ( M_{n-1} M_{n-2} \cdots M_{0} )_{b} \\
    \overline{r} = ( \overline{r}_{n-1} \overline{r}_{n-2} \cdots \overline{r}_{0} )_{b}
\end{gather*}

Each of the values in
\begin{math}Z_{n-1} \cdots Z_{0}\end{math},
\begin{math}X_{n-1} \cdots X_{0}\end{math},
\begin{math}Y_{n-1} \cdots Y_{0}\end{math},
\begin{math}M_{n-1} \cdots M_{0}\end{math},
\begin{math}\overline{r}_{n-1}\cdots \overline{r}_{0}\end{math} represents one base-\begin{math}b\end{math} digit (a 32-bit word).

\begin{math}Z_{n-1}\end{math},
\begin{math}X_{n-1}\end{math},
\begin{math}Y_{n-1}\end{math},
\begin{math}M_{n-1}\end{math} and
\begin{math}\overline{r}_{n-1}\end{math} are the most significant bits of
\begin{math}Z\end{math},
\begin{math}X\end{math},
\begin{math}Y\end{math},
\begin{math}M\end{math}, and \begin{math}\overline{r}\end{math}, while
\begin{math}Z_{0}\end{math},
\begin{math}X_{0}\end{math},
\begin{math}Y_{0}\end{math},
\begin{math}M_{0}\end{math} and
\begin{math}\overline{r}_{0}\end{math} are the least significant bits.

If we define \begin{math}R = b^{n}\end{math}, the additional argument $\overline{r}$ can be calculated as \begin{math}\overline{r} = R^{2} \bmod M\end{math}.

Also, argument $M^{\prime}$ can be calculated using the formula below:
\[
 M^{\prime} = - M^{-1} \bmod b
\]
where, $M^{-1}$ is the \href{https://en.wikipedia.org/wiki/Modular\_multiplicative\_inverse}{modular multiplicative inverse} of $M$, and it can be calculated with the extended binary GCD algorithm.

Large-number modular exponentiation on the \chipname{} can be implemented as follows:

\begin{enumerate}

    \item Write 1 or 0 to the \hyperref[fielddesc:RSAINTENA]{RSA\_INT\_ENA} field of the register \hyperref[regdesc:RSAINTENAREG]{RSA\_INT\_ENA\_REG} to enable or disable the interrupt function.

    \item Configure relevant registers:

    \begin{enumerate}
        \item Write \begin{math}( \frac{N}{32} - 1 )\end{math} to the \hyperref[regdesc:RSAMODEREG]{RSA\_MODE\_REG} register.

        \item Write \begin{math}M^{\prime}\end{math} to the \hyperref[regdesc:RSAMPRIMEREG]{RSA\_M\_PRIME\_REG} register.

        \item Configure registers related to the acceleration options, which are described later in Section \ref{sec:Accelerate options}.

    \end{enumerate}

    \item Write
        \begin{math}X_{i}\end{math},
        \begin{math}Y_{i}\end{math},
        \begin{math}M_{i}\end{math} and
        \begin{math}\overline{r}_{i}\end{math} for \begin{math}i \in \{ 0, 1, \dots, n-1 \} \end{math}
        to memory blocks \hyperref[tab:rsa-mem]{RSA\_X\_MEM}, \hyperref[tab:rsa-mem]{RSA\_Y\_MEM}, \hyperref[tab:rsa-mem]{RSA\_M\_MEM} and \hyperref[tab:rsa-mem]{RSA\_Z\_MEM}. The capacity of each memory block is 96 words. Each word of each memory block can store one base-\begin{math}b\end{math} digit. The memory blocks use the little endian format for storage, i.e., the least significant digit of each number is in the lowest address.

       Users need to write data to each memory block only according to the length of the number; data beyond this length is ignored.

    \item Write 1 to the \hyperref[fielddesc:RSASETSTARTMODEXP]{RSA\_SET\_START\_MODEXP} field of the register \hyperref[regdesc:RSASETSTARTMODEXPREG]{RSA\_SET\_START\_MODEXP\_REG} to start computation.

    \item Wait for the completion of computation, which happens when the content of \hyperref[fielddesc:RSAQUERYIDLE]{RSA\_QUERY\_IDLE} becomes 1 or the RSA interrupt occurs.

    \item Read the result \begin{math}Z_{i}\end{math}
        for \begin{math}i \in \{ 0, 1, \dots, n-1 \} \end{math}
        from \hyperref[tab:rsa-mem]{RSA\_Z\_MEM}.

    \item Write 1 to \hyperref[fielddesc:RSACLEARINTERRUPT]{RSA\_CLEAR\_INTERRUPT} to clear the interrupt, if you have the interrupt enabled.

\end{enumerate}

After the computation, the \hyperref[regdesc:RSAMODEREG]{RSA\_MODE\_REG} register, memory blocks \hyperref[tab:rsa-mem]{RSA\_Y\_MEM} and \hyperref[tab:rsa-mem]{RSA\_M\_MEM}, as well as the \hyperref[regdesc:RSAMPRIMEREG]{RSA\_M\_PRIME\_REG} remain unchanged. However, \begin{math}X_{i}\end{math} in \hyperref[tab:rsa-mem]{RSA\_X\_MEM} and \begin{math}\overline{r}_{i}\end{math} in \hyperref[tab:rsa-mem]{RSA\_Z\_MEM} computation are overwritten, and only these overwritten memory blocks need to be re-initialized before starting another computation.

%%%%%%%%%%%%%%%%%%%%%%%%%%%%%%%%%%%%%%%%%%%%%%%%%
%%%%%%%%%%%%%%%%%%%%%%%%%%%%%%%%%%%%%%%%%%%%%%%%%
\subsection{Large-Number Modular Multiplication}

Large-number modular multiplication performs \begin{math}Z = X \times Y \bmod M\end{math}. This computation is based on Montgomery multiplication.
Therefore, similar to the large-number modular exponentiation, two additional arguments are needed -- \begin{math}\overline{r}\end{math} and
\begin{math}M^{\prime}\end{math}, which need to be calculated in advance by software.

The RSA accelerator supports large-number modular multiplication with operands of 96 different lengths.

The computation can be executed as follows:

\begin{enumerate}

    \item Write 1 or 0 to the \hyperref[regdesc:RSAINTENAREG]{RSA\_INT\_ENA\_REG} register to enable or disable the interrupt function.

    \item Configure relevant registers:

    \begin{enumerate}
        \item Write \begin{math}( \frac{N}{32} - 1 )\end{math} to the \hyperref[regdesc:RSAMODEREG]{RSA\_MODE\_REG} register.

        \item Write \begin{math}M^{\prime}\end{math} to the \hyperref[regdesc:RSAMPRIMEREG]{RSA\_M\_PRIME\_REG} register.

    \end{enumerate}

    \item Write
        \begin{math}X_{i}\end{math},
        \begin{math}Y_{i}\end{math},
        \begin{math}M_{i}\end{math}, and
        \begin{math}\overline{r}_{i}\end{math} for \begin{math}i \in \{ 0, 1, \dots, n-1 \} \end{math} to memory blocks \hyperref[tab:rsa-mem]{RSA\_X\_MEM}, \hyperref[tab:rsa-mem]{RSA\_Y\_MEM}, \hyperref[tab:rsa-mem]{RSA\_M\_MEM}, and \hyperref[tab:rsa-mem]{RSA\_Z\_MEM}, respectively. The capacity of each memory block is 96 words. Each word of each memory block can store one base-\begin{math}b\end{math} digit. The memory blocks use the little endian format for storage, i.e., the least significant digit of each number is in the lowest address.

        Users need to write data to each memory block only according to the length of the number; data beyond this length are ignored.

    \item Write 1 to the \hyperref[fielddesc:RSASETSTARTMODMULT]{RSA\_SET\_START\_MODMULT} field.

    \item Wait for the completion of computation, which happens when the content of \hyperref[fielddesc:RSAQUERYIDLE]{RSA\_QUERY\_IDLE} becomes 1 or the RSA interrupt occurs.

    \item Read the result \begin{math}Z_{i}\end{math} for \begin{math}i \in \{ 0, 1, \dots, n-1 \} \end{math} from \hyperref[tab:rsa-mem]{RSA\_Z\_MEM}.
    \item Write 1 to \hyperref[fielddesc:RSACLEARINTERRUPT]{RSA\_CLEAR\_INTERRUPT} to clear the interrupt, if you have the interrupt enabled.

\end{enumerate}

After the computation, the length of operands in \hyperref[regdesc:RSAMODEREG]{RSA\_MODE\_REG}, the \begin{math}X_{i}\end{math} in memory \hyperref[tab:rsa-mem]{RSA\_X\_MEM}, the \begin{math}Y_{i}\end{math} in memory \hyperref[tab:rsa-mem]{RSA\_Y\_MEM}, the \begin{math}M_{i}\end{math} in memory \hyperref[tab:rsa-mem]{RSA\_M\_MEM}, and the \begin{math}M^{\prime}\end{math} in memory \hyperref[regdesc:RSAMPRIMEREG]{RSA\_M\_PRIME\_REG} remain unchanged. However, the \begin{math}\overline{r}_{i}\end{math} in memory \hyperref[tab:rsa-mem]{RSA\_Z\_MEM} has already been overwritten, and only this overwritten memory block needs to be re-initialized before starting another computation.

%%%%%%%%%%%%%%%%%%%%%%%%%%%%%%%%%%%%%%%%%%%%%%%%%
%%%%%%%%%%%%%%%%%%%%%%%%%%%%%%%%%%%%%%%%%%%%%%%%%
\subsection{Large-Number Multiplication}

Large-number multiplication performs \begin{math}Z = X \times Y\end{math}. The length of result \begin{math}Z\end{math} is twice that of operand \begin{math}X\end{math} and operand \begin{math}Y\end{math}. Therefore, the RSA accelerator only supports large-number multiplication with operand length \begin{math}N = 32 \times x\end{math}, where \begin{math}x \in \{ 1, 2, 3, \dots, 48\}\end{math}. The length \begin{math}\hat{N}\end{math} of result \begin{math}Z\end{math} is \begin{math}2 \times N\end{math}.

The computation can be executed as follows:

\begin{enumerate}

    \item Write 1 or 0 to the \hyperref[regdesc:RSAINTENAREG]{RSA\_INT\_ENA\_REG} register to enable or disable the interrupt function.

    \item Write \begin{math}( \frac{\hat{N}}{32} - 1 )\end{math}, i.e., \begin{math}( \frac{N}{16} - 1 )\end{math} to the \hyperref[regdesc:RSAMODEREG]{RSA\_MODE\_REG} register.

    \item Write \begin{math}{X}_{i}\end{math} and \begin{math}{Y}_{i}\end{math} for \begin{math} \in \{ 0, 1, \dots, n-1 \} \end{math} to memory blocks \hyperref[tab:rsa-mem]{RSA\_X\_MEM} and \hyperref[tab:rsa-mem]{RSA\_Z\_MEM}. Each word of each memory block can store one base-\begin{math}b\end{math} digit. The memory blocks use the little endian format for storage, i.e., the least significant digit of each number is in the lowest address. \begin{math}n\end{math} is \begin{math} \frac{N}{32} \end{math}.

         Write \begin{math}{X}_{i}\end{math} for \begin{math}i \in \{ 0, 1, \dots, n-1 \} \end{math} to the address of the \begin{math}i\end{math} words of the \hyperref[tab:rsa-mem]{RSA\_X\_MEM} memory block. Note that \begin{math}{Y}_{i}\end{math} for \begin{math}i \in \{ 0, 1, \dots, n-1 \} \end{math} will not be written to the address of the \begin{math}i\end{math} words of the \hyperref[tab:rsa-mem]{RSA\_Z\_MEM} register, but the address of the \begin{math}n+i\end{math} words, i.e., the base address of the \hyperref[tab:rsa-mem]{RSA\_Z\_MEM} memory plus the address offset \begin{math}4 \times (n+i) \end{math}.

        Users need to write data to each memory block only according to the length of the number; data beyond this length is ignored.

    \item Write 1 to the \hyperref[fielddesc:RSASETSTARTMULT]{RSA\_SET\_START\_MULT} register.

    \item Wait for the completion of computation, which happens when the content of \hyperref[fielddesc:RSAQUERYIDLE]{RSA\_QUERY\_IDLE} becomes 1 or the RSA interrupt occurs.

    \item Read the result \begin{math}Z_{i}\end{math}
          for \begin{math}i \in \{ 0, 1, \dots, \hat{n}-1 \} \end{math} from the \hyperref[tab:rsa-mem]{RSA\_Z\_MEM} register.
        \begin{math}\hat{n}\end{math} is \begin{math} 2 \times n \end{math}.

    \item Write 1 to \hyperref[fielddesc:RSACLEARINTERRUPT]{RSA\_CLEAR\_INTERRUPT} to clear the interrupt, if you have the interrupt enabled.

\end{enumerate}

After the computation, the length of operands in \hyperref[regdesc:RSAMODEREG]{RSA\_MODE\_REG} and the \begin{math}X_{i}\end{math} in memory \hyperref[tab:rsa-mem]{RSA\_X\_MEM} remain unchanged. However, the \begin{math}{Y}_{i}\end{math} in memory \hyperref[tab:rsa-mem]{RSA\_Z\_MEM} has already been overwritten, and only this overwritten memory block needs to be re-initialized before starting another computation.
%%%%%%%%%%%%%%%%%%%%%%%%%%%%%%%%%%%%%%%%%%%%%%%%%
%%%%%%%%%%%%%%%%%%%%%%%%%%%%%%%%%%%%%%%%%%%%%%%%%
\subsection{Options for Additional Acceleration} \label{sec:Accelerate options}
The \chipname{} RSA accelerator also provides \hyperref[regdesc:RSASEARCHENABLEREG]{SEARCH} and \hyperref[regdesc:RSACONSTANTTIMEREG]{CONSTANT\_TIME} options that can be configured to further accelerate the large-number modular exponentiation. By default, both options are configured as no additional acceleration.

Users can choose to use one or two of these options to further accelerate the computation. Note that, even when none of these two options is configured, using the hardware RSA accelerator is still much faster than implementing the RSA algorithm in software.

To be more specific, when neither of these two options are configured for additional acceleration, the time required to calculate \begin{math}Z = X^{Y} \bmod M\end{math} is solely determined by the lengths of operands.  When either or both of these two options are configured for additional acceleration, the time required is also correlated with the 0/1 distribution of \begin{math}Y\end{math}.

To better illustrate how these two options work, first assume \begin{math}Y \end{math} is represented in binaries as

\begin{gather*}
Y = ( \widetilde{Y}_{N-1} \widetilde{Y}_{N-2} \cdots \widetilde{Y}_{t+1} \widetilde{Y}_{t} \widetilde{Y}_{t-1} \cdots \widetilde{Y}_{0} )_{2} \\
\end{gather*}
where,
\begin{itemize}
    \item \begin{math} N \end{math} is the length of \begin{math}Y \end{math},
%%%%
    \item \begin{math}\widetilde{Y}_{t}\end{math} is 1,
%%%%
    \item \begin{math}\widetilde{Y}_{N-1}\end{math}, \begin{math}\widetilde{Y}_{N-2}\end{math}, \dots, \begin{math}\widetilde{Y}_{t+1}\end{math} are all equal to 0,
%%%%
    \item and \begin{math}\widetilde{Y}_{t-1}\end{math}, \begin{math}\widetilde{Y}_{t-2}\end{math}, \dots, \begin{math}\widetilde{Y}_{0}\end{math} are either 0 or 1 but exactly \begin{math}m\end{math} bits should be equal to 0 and \begin{math}t -m\end{math} bits 1, i.e., the Hamming weight of \begin{math}\widetilde{Y}_{t-1} \widetilde{Y}_{t-2}, \cdots, \widetilde{Y}_{0}\end{math} is \begin{math}t -m\end{math}.
\end{itemize}

When either of these two options is configured for additional acceleration:

\begin{itemize}
 \item SEARCH Option (Configuring \hyperref[regdesc:RSASEARCHENABLEREG]{RSA\_SEARCH\_ENABLE} to 1 for additional acceleration)
    \begin{itemize}
        \item The accelerator ignores the bit positions of \begin{math} \widetilde{Y}_{i} \end{math}, where \begin{math} i > \alpha \end{math}. Search position \begin{math} \alpha \end{math} is set by configuring the \hyperref[regdesc:RSASEARCHPOSREG]{RSA\_SEARCH\_POS\_REG} register. Set \begin{math} \alpha \end{math} to a number smaller than \begin{math}N\end{math}-1, which otherwise leads to the same result as if this option is not used for additional acceleration. The best acceleration performance can be achieved by setting \begin{math} \alpha \end{math} to \begin{math} t \end{math}, in which case all the 0 bits in \begin{math}\widetilde{Y}_{N-1}\end{math}, \begin{math}\widetilde{Y}_{N-2}\end{math}, \dots, \begin{math}\widetilde{Y}_{t+1}\end{math} are ignored during the calculation. Note that if you set \begin{math} \alpha \end{math} to be less than \begin{math} t \end{math}, then the result of the modular exponentiation \begin{math}Z = X^{Y} \bmod M \end{math} will be incorrect.%
    \item Note that this option compromises the security because it ignores some bits, which essentially shortens the key length, thus should not be enabled for applications with high security requirement.
    \end{itemize}

\item CONSTANT\_TIME Option (Configuring \hyperref[regdesc:RSACONSTANTTIMEREG]{RSA\_CONSTANT\_TIME\_REG} to 0 for additional acceleration)
    \begin{itemize}
        \item The accelerator speeds up the calculation by simplifying the calculation concerning the 0 bits of \begin{math} Y\end{math}. Therefore, the higher the proportion of bits 0 against bits 1, the better is the acceleration performance.
        \item Note that this option also compromises the security because its time cost correlates with the 0/1 distribution of the key, which may be used in a Side Channel Attack (SCA), thus should not be enabled for applications with high security requirement.
    \end{itemize}
\end{itemize}

Below is an example to demonstrate the performance of the RSA accelerator under different combinations of \hyperref[regdesc:RSASEARCHENABLEREG]{SEARCH} and \hyperref[regdesc:RSACONSTANTTIMEREG]{CONSTANT\_TIME} configuration. \\
In this example:
\begin{itemize}
    \item We perform \begin{math}Z = X^{Y} \bmod M\end{math}
    \item \begin{math}N\end{math} = 3072
    \item \begin{math}Y\end{math} = 65537
    \item \begin{math}X\end{math} and \begin{math}M\end{math} are taken at random
    \item When the SEARCH option is enabled, \begin{math} \alpha \end{math}, i.e., the position register \hyperref[regdesc:RSASEARCHPOSREG]{RSA\_SEARCH\_POS\_REG}, is set to 16.
\end{itemize}

Table \ref{tab:accelerationeffect} below demonstrates the time cost in clock cycles under different combinations of \hyperref[regdesc:RSASEARCHENABLEREG]{SEARCH} and \hyperref[regdesc:RSACONSTANTTIMEREG]{CONSTANT\_TIME} configuration when performing \begin{math}Z = X^{Y} \bmod M\end{math} as described above.

\begin{longtable}{ | l | l | r | }
    \caption{Acceleration Performance}\label{tab:accelerationeffect}
    \\\hline
    \rowcolor{lightgray}
        \textbf{SEARCH Option} & \textbf{CONSTANT\_TIME Option} & \textbf{Time Cost (clock cycle)}  \\ \hline
        \endfirsthead
        \multicolumn{3}{c}{\bfseries \tablename\ \thetable{} -- cont'd from previous page}\\\hline

        \rowcolor{lightgray}
        \textbf{SEARCH Option} & \textbf{CONSTANT\_TIME Option} & \textbf{Time Cost (clock cycle)}  \\ \hline
        \endhead

        \multicolumn{3}{r}{\textbf{Cont'd on next page}} \\
        \endfoot
        \endlastfoot
        \hline
    No acceleration & No acceleration   & $174.7 \times 10^6$   \\ \hline
    Acceleration     & No acceleration   & $1.023 \times 10^6$   \\ \hline
    No acceleration & Acceleration      & $0.546 \times 10^6$   \\ \hline
    Acceleration    & Acceleration      & $0.540 \times 10^6$   \\ \hline
\end{longtable}

As shown in Table \ref{tab:accelerationeffect}:

\begin{itemize}
    \item The time cost is the biggest when none of these two options is configured for additional acceleration.
    \item The time cost is the smallest when both of these two options are configured for additional acceleration.
    \item The time cost can be dramatically reduced when either or both option(s) are configured for additional acceleration.
\end{itemize}

%%%%%%%%%%%%%%%%%%%%%
%%% Interrupts %%%
%%%%%%%%%%%%%%%%%%%%%

\section{Interrupts}

\chipname{}'s RSA accelerator can generate the following interrupt signal that will be sent to the \hyperref[mod:intmtrx]{Interrupt Matrix}.

    \begin{itemize}
        \item RSA\_INTR
    \end{itemize}

There is one internal interrupt source from the RSA accelerator that can generate the above interrupt signal. The interrupt source from the RSA accelerator is listed with its trigger condition and the resulting interrupt signal in Table \ref{tab:rsa-int-sources}.

\begin{longtable}{ | L{5cm} | L{6.5cm} | L{4cm} | }
\caption{RSA's Internal Interrupt Source} \label{tab:rsa-int-sources}
\\ \hline
\rowcolor{lightgray}
\textbf{Internal Interrupt Source} & \textbf{Trigger Condition}  &  \textbf{Interrupt Signal}\\ \hline
\label{int:RSACALCDONEINT}RSA\_CALC\_DONE\_INT  & Completion of an RSA calculation &  RSA\_INTR  \\ \hline
\end{longtable}

\begin{tiplisting}
For definitions of \textit{interrupt}, \textit{interrupt signal}, \textit{interrupt source}, and their correlations, please refer to Chapter \ref{mod:intmtrx} \textit{\nameref{mod:intmtrx}} > Section \ref{sec:intmtx-int-terminology} \textit{\nameref{sec:intmtx-int-terminology}}.
\end{tiplisting}

%%%%%%%%%%%%%%%%%%%%%%%%%%%%%%%%%%%%%%%%%%%%%%%%%%%%%%%%%%%%%%%%%%%%%%%%%%%%%%%%%
%%%%%%%%%%%%%%%%%%%%%%%%%%%%%%%%%%%%%%%%%%%%%%%%%%%%%%%%%%%%%%%%%%%%%%%%%%%%%%%%%
\newpage
\section{Memory Summary}\label{sec:mem-sum}
The addresses in this section are relative to the RSA accelerator base address provided in Table \ref{tab:sysmem-base-address} in Chapter \ref{mod:sysmem} \textit{\nameref{mod:sysmem}}.

The abbreviations given in Column \textbf{Access} are explained in Section \hyperref[glossary-access-types]{\textit{Access Types for Registers}}.

\subfile{\modulefiles/25-RSA-mem-summ__EN}


%%%%%%%%%%%%%%%%%%%%%%%%
%%% Register Summary %%%
%%%%%%%%%%%%%%%%%%%%%%%%


\hypertarget{rsa-reg-summ}{}
\section{Register Summary}
\label{sec:rsa-reg-summ}


The addresses in this section are relative to the RSA accelerator base address provided in Table \ref{tab:sysmem-base-address} in Chapter \ref{mod:sysmem} \textit{\nameref{mod:sysmem}}.

The abbreviations given in Column \textbf{Access} are explained in Section \hyperref[glossary-access-types]{\textit{Access Types for Registers}}.

\subfile{\modulefiles/25-RSA-reg-summ__EN}


%%%%%%%%%%%%%%%%%
%%% Registers %%%
%%%%%%%%%%%%%%%%%

\newpage
\section{Registers}

The addresses in this section are relative to the RSA accelerator base address provided in Table \ref{tab:sysmem-base-address} in Chapter \ref{mod:sysmem} \textit{\nameref{mod:sysmem}}.

For how to program reserved fields, please refer to Section \hyperref[programming-reserved-field]{Programming Reserved Register Field}.


\subfile{\modulefiles/25-RSA-reg__EN}

\end{document}
